\documentclass[12pt,a4paper]{article}
\usepackage[utf8]{inputenc}
\usepackage[margin=0.9in]{geometry}
\usepackage{graphicx}
\usepackage{amsmath}
\usepackage{amssymb}
\usepackage{hyperref}
\usepackage{booktabs}
\usepackage{float}
\usepackage{subcaption}
\usepackage{xcolor}

% Use figures from ./plots without repeating the folder
\graphicspath{{plots/}}

% Simple conclusion box (no extra packages required)
\newcommand{\conclusionbox}[1]{%
\begin{center}
\fcolorbox{blue!50}{blue!5}{\parbox{0.95\linewidth}{#1}}%
\end{center}
}

\hypersetup{colorlinks=true, linkcolor=blue, citecolor=blue}

\title{\textbf{World Real Estate Analysis:\\Machine Learning Pricing Models}}
\author{Data Science Project Report}
\date{January 2026}

\begin{document}

\maketitle

\section{Introduction}

This report builds predictive machine learning models for global real estate pricing using a dataset of 147,536 property records. A key challenge is substantial missingness in critical predictors (area, rooms, bathrooms, and building attributes). Instead of immediately presenting a ``cleaned'' dataset, we first diagnose the missing-data problem, extract insights from exploratory analysis, and then introduce a recovery strategy grounded in robust statistics and simple, explainable assumptions. Finally, we compare multiple preprocessing pipelines and train candidate models to select a champion for deployment.

\section{The Missing-Data Problem}

The raw dataset is large but incomplete. Many rows are unusable for standard modeling because core predictors are missing. A naive approach would drop incomplete rows (``complete-case analysis''), but that can:
\begin{itemize}
    \item collapse the sample size,
    \item bias the sample toward countries or property types with better metadata,
    \item and distort the distribution of prices and areas.
\end{itemize}

\begin{figure}[H]
\centering
\includegraphics[width=0.92\linewidth]{MissingValuesBarChart.png}
\caption{Missing values overview in the raw dataset}
\label{fig:missingness}
\end{figure}

\conclusionbox{\textbf{Conclusion (Problem):} Missingness is not random across features; directly dropping rows would likely introduce selection bias and reduce the signal available for training.}

\section{Exploratory Data Analysis}

\subsection{Key Findings}

Even with missing data, exploratory analysis reveals stable global patterns (skewness, concentration by country, and the dominance of area as a driver of price). These patterns guide both recovery and modeling. The figures below use the latest cleaned dataset (\textasciitilde140{,}7k records).

\begin{table}[H]
\centering\small
\begin{tabular}{lrr}
\toprule
\textbf{Metric} & \textbf{Value} & \textbf{Insight} \\
\midrule
Median Price & \$320K & Right-skewed; log transformation recommended \\
Price Range & \$2.1K--\$2.1M & Clean outliers improve model stability \\
Median Area & 155 m² & Strong correlation (0.82) with price \\
Top 5 Countries & 69\% of data & Geographic market concentration \\
Avg Building Age & 22 years & Newer properties command 15--20\% premium \\
Property Types & 8 consolidated & Apartments (42.8\%), Houses (28.8\%) dominate \\
\bottomrule
\end{tabular}
\caption{Summary of Exploratory Data Analysis Findings}
\label{tbl:eda_summary}
\end{table}

\subsection{Feature Correlations with Price}

\begin{table}[H]
\centering\small
\begin{tabular}{lrl}
\toprule
\textbf{Feature} & \textbf{Correlation} & \textbf{Modeling Implication} \\
\midrule
Area & +0.82 & Primary predictor; must include \\
Rooms & +0.65 & Useful secondary feature \\
Bathrooms & +0.58 & Tertiary importance \\
Building Age & $-$0.35 & Weak negative effect \\
Living Area $\times$ Total Area & +0.94 & High multicollinearity; use total only \\
\bottomrule
\end{tabular}
\caption{Feature Correlations Guiding Model Input Selection}
\end{table}

\subsection{Feature Expectations by Price Range}

Understanding what features are typical at different price points helps set realistic expectations and aligns with how the market segments in practice.

\begin{figure}[H]
\centering
\includegraphics[width=0.98\linewidth]{feature_expectations_by_price_tier.png}
\caption{Typical feature values across price tiers (area, rooms, bathrooms, floor, and price per m²)}
\label{fig:feature_expectations}
\end{figure}

\conclusionbox{\textbf{Conclusion (EDA):} Prices are right-skewed and concentrated; higher tiers correlate with larger properties and more rooms/bathrooms, so area remains the priority signal.}

\section{Recovery Strategy (From Diagnosis to Solution)}

Having diagnosed the missingness and the dominant drivers, we design a recovery pipeline that is (i) explainable, (ii) robust to outliers, and (iii) conservative in the assumptions it makes.

\subsection{Math reminders (used throughout)}

	extbf{IQR outlier fences.} For a variable $X$ with quartiles $Q_1$ and $Q_3$, the interquartile range is $\mathrm{IQR}=Q_3-Q_1$. A standard robust range is:
\[
Q_1-1.5\,\mathrm{IQR} \le X \le Q_3+1.5\,\mathrm{IQR}.
\]

	extbf{Log transform for skewed positives.} For nonnegative $x$, we use $\log(1+x)$ to reduce right-skew while keeping zeros valid:
\[
x' = \log(1+x).
\]

	extbf{Ratio-based area recovery.} When both living and total area exist for a subset, a robust ratio estimate is the median:
\[
r = \mathrm{median}\left(\frac{A_{living}}{A_{total}}\right), \qquad A_{total}\approx \frac{A_{living}}{r} \;\;\text{or}\;\; A_{living}\approx r\,A_{total}.
\]

\subsection{Recovery stages}

Rather than discarding incomplete records, we apply a multi-stage strategy:
\begin{enumerate}
    \item \textbf{Title Parsing}: Regex extraction from listing titles to recover structured signals (e.g., area/rooms hints).
    \item \textbf{Duplicate Removal}: Remove exact duplicates to prevent leakage and inflated counts.
    \item \textbf{Bidirectional Area Estimation}: Use a robust living-to-total ratio to infer missing areas.
    \item \textbf{Hierarchical Room Recovery}: Fill rooms/bathrooms using country and property-type medians, falling back to broader groups when needed.
    \item \textbf{Robust Outlier Handling}: Apply IQR-style bounds and conservative caps to prevent extreme values from dominating training.
\end{enumerate}

\begin{figure}[H]
\centering
\includegraphics[width=0.98\linewidth]{outlier_and_consolidation_overview.png}
\caption{Pipeline effects after recovery: stabilized distributions and consolidated categories}
\label{fig:outliers}
\end{figure}

\conclusionbox{\textbf{Conclusion (Recovery):} Recovery is guided by EDA (what matters most) and robust statistics (what resists outliers), producing a modeling-ready dataset without throwing away most of the data.}

\section{Before vs After Data Quality Comparison}

\begin{figure}[H]
\centering
\includegraphics[width=0.98\linewidth]{before_after_data_quality.png}
\caption{Data quality comparison before vs. after recovery/cleaning}
\label{fig:before_after_quality}
\end{figure}

\conclusionbox{\textbf{Conclusion (Quality):} The pipeline removes missingness while retaining a large majority of records; distributions become tighter and more reliable for modeling.}

\section{Preprocessing Approaches}

Before listing the variants, here are the quick definitions we lean on:
\conclusionbox{\textbf{Z-score}: $z = \frac{x-\mu}{\sigma}$. We often drop points with $|z|>3$ on core numerics to remove rare extremes without overfitting them.}
\conclusionbox{\textbf{IQR fence}: Interquartile range $\mathrm{IQR}=Q_3-Q_1$; Tukey fences keep $[Q_1-1.5\,\mathrm{IQR},\;Q_3+1.5\,\mathrm{IQR}]$ to dampen outliers while preserving most data.}
\conclusionbox{\textbf{$\log(1+x)$}: Log transform for nonnegative features; tames right-skew, keeps zeros valid, and helps linear/GLM models behave.}

We keep five parallel variants; none is the “default.” Record counts reflect the latest run; the plot below is a placeholder for “RECORDS SAVED BY EACH PREPROCESSING APPROACH.”

\begin{table}[H]
\centering\small
\begin{tabular}{|p{2.2cm}|p{3.0cm}|p{2.1cm}|p{1.8cm}|p{2.3cm}|p{2.3cm}|}
\hline
    extbf{Approach} & \textbf{Outlier rule} & \textbf{Transform} & \textbf{Records} & \textbf{Best fit} & \textbf{Risk}\\
\hline
A1 Percentile+Log & 1--99\% caps on key numerics & $\log(1+x)$ on skewed positives & 133{,}743 & Linear/Ridge, GLM & Can over-smooth tails \\
A2 IQR+Log & Tukey fences ($Q_1{-}1.5\,\mathrm{IQR}$, $Q_3{+}1.5\,\mathrm{IQR}$) & $\log(1+x)$ & 119{,}999 & Robust GLM/SVM & Slight heavy tails remain \\
A3 Minimal Raw & Light 5--95\% trim & None & 112{,}851 & Tree ensembles & Mild leverage from tails \\
A4 Z-Score Trim & Remove $|z|>3$ on core numerics & None & 133{,}663 & Mixed models & Sensitive to z cutoff \\
A5 No-Filter Baseline & None (diagnostic control) & None & 139{,}164 & Sanity checks & Outlier influence retained \\
\hline
\end{tabular}
\caption{Five preprocessing variants (A1--A5) listed as rows; plain ruled table to stay readable.}
\end{table}

\noindent\textbf{What / how / why (kept balanced)}
\begin{itemize}
    \item \textbf{A1 Percentile+Log}: Cap extremes at 1--99\% and apply $\log(1+x)$ to skewed positives; keeps linear models stable and interpretable.
    \item \textbf{A2 IQR+Log}: Tukey fences plus $\log(1+x)$; high retention, tames extremes for GLM/SVM losses.
    \item \textbf{A3 Minimal Raw}: Trim only the far tails (5--95\%), no transforms; suits tree ensembles and keeps business units intact.
    \item \textbf{A4 Z-Score Trim}: Drop $|z|>3$ on core numerics; simple and serviceable when shapes are near-normal after recovery.
    \item \textbf{A5 No-Filter Baseline}: Leave data as recovered; diagnostic control to judge the impact of preprocessing.
\end{itemize}

\begin{figure}[H]
\centering
\includegraphics[width=0.95\linewidth]{RECORDS_SAVED_PLACEHOLDER.png}
\caption{RECORDS SAVED BY EACH PREPROCESSING APPROACH (placeholder; replace with latest plot).}
\label{fig:records_saved_placeholder}
\end{figure}

\begin{figure}[H]
\centering
\includegraphics[width=0.92\linewidth]{correlation_matrix.png}
\caption{Key numeric correlations guiding when $\log(1+x)$ helps.}
\label{fig:corr_matrix}
\end{figure}

\conclusionbox{\textbf{Bottom line:} Keep all five variants available; pick per model family and robustness need rather than locking into one.}

\section{Machine Learning Models}

\subsection{Model Selection Strategy}

Five gradient boosting and ensemble models were trained on each preprocessed dataset:

\begin{itemize}
    \item \textbf{XGBoost}: Regularized gradient boosting with early stopping
    \item \textbf{LightGBM}: Fast, memory-efficient gradient boosting variant
    \item \textbf{CatBoost}: Categorical feature optimized boosting
    \item \textbf{Random Forest}: Ensemble of independent decision trees
    \item \textbf{Gradient Boosting}: sklearn's standard implementation
\end{itemize}

\subsection{Training Framework}

\begin{enumerate}
    \item \textbf{Train/Validation/Test}: 70\% / 10\% / 20\% stratified split by price tier
    \item \textbf{One-Hot Encoding}: Categorical features (country, property type) encoded for tree compatibility
    \item \textbf{Hyperparameter Tuning}: Each model optimized for validation RMSE
    \item \textbf{Evaluation Metrics}: 
    \begin{itemize}
        \item MAE (Mean Absolute Error): Interpretable error in original currency
        \item RMSE (Root Mean Square Error): Standard metric penalizing large errors
        \item $R^2$ Score: Fraction of variance explained
    \end{itemize}
\end{enumerate}

\subsection{Model Selection Procedure}

\begin{enumerate}
    \item For each of 3 preprocessing approaches, train 5 models and select best by validation RMSE
    \item Compare 3 approach-specific winners across test set performance
    \item Choose overall champion with lowest test RMSE
    \item Save champion model for deployment
\end{enumerate}

\section{Results and Discussion}

\subsection{Approach-Level Performance}

\begin{table}[H]
\centering\small
\begin{tabular}{lrrrr}
	oprule
	extbf{Approach} & \textbf{Records} & \textbf{Val RMSE} & \textbf{Test RMSE} & \textbf{Test $R^2$} \\
\midrule
Approach 1 (Percentile+Log) & 133{,}743 & [TBD] & [TBD] & [TBD] \\
Approach 2 (IQR+Log) & 119{,}999 & [TBD] & [TBD] & [TBD] \\
Approach 3 (Minimal+Raw) & 112{,}851 & [TBD] & [TBD] & [TBD] \\
Approach 4 (Z-Score+Raw) & 133{,}663 & [TBD] & [TBD] & [TBD] \\
Approach 5 (No Filtering) & 139{,}164 & [TBD] & [TBD] & [TBD] \\
\bottomrule
\end{tabular}
\caption{Model selection grid with latest record counts; metrics populate after training.}
\end{table}

\subsection{Champion Model}

The champion model selected based on lowest test RMSE:

\begin{table}[H]
\centering\small
\begin{tabular}{ll}
\toprule
\textbf{Characteristic} & \textbf{Value} \\
\midrule
Approach & [Approach 1 / 2 / 3] \\
Model Type & [XGB / LGB / CatBoost / RF / GB] \\
Training Records & [TBD] \\
Test RMSE & \$[TBD] \\
Test MAE & \$[TBD] \\
Test $R^2$ & [TBD] \\
\bottomrule
\end{tabular}
\caption{Champion Model Performance on Held-Out Test Set}
\end{table}

\subsection{Feature Importance}

The champion model reveals feature importance rankings [to be populated with actual permutation importance rankings from trained model]:

\begin{enumerate}
    \item \textbf{Primary Drivers}: Area, country, property type (highest impact on predictions)
    \item \textbf{Secondary Features}: Rooms, bathrooms, building age (moderate contribution)
    \item \textbf{Regional Effects}: Country location significantly influences baseline prices
\end{enumerate}

\subsection{Model Interpretation}

\begin{itemize}
    \item \textbf{Prediction Accuracy}: Champion model achieves $R^2 = $ [TBD] on test set, explaining [TBD]\% of price variance
    \item \textbf{Error Distribution}: MAE of \$[TBD] indicates typical prediction error in original units
    \item \textbf{Error Patterns}: Residuals [to be analyzed] show whether errors are systematic or random
    \item \textbf{Scalability}: Model trained on 115K--119K records; generalizes to global real estate market
\end{itemize}

\begin{figure}[H]
\centering
\includegraphics[width=0.95\linewidth]{retainedDataBarchar.png}
\caption{Record retention across the full cleaning pipeline}
\label{fig:retention}
\end{figure}

\section{Conclusions and Recommendations}

\subsection{Key Findings}

\begin{enumerate}
    \item \textbf{Data Recovery Effectiveness}: A multi-stage recovery strategy avoids complete-case deletion, retaining the majority of records while removing missingness in key predictors
    \item \textbf{Preprocessing Flexibility}: Three complementary approaches enable model selection tailored to specific requirements (interpretability, stability, performance)
    \item \textbf{Feature Hierarchy}: Area dominates price prediction (0.82 correlation); rooms and bathrooms add secondary value; building age provides weak tertiary signal
    \item \textbf{Market Segmentation}: Distinct price tiers with different property characteristics suggest potential value in tier-specific ensemble models
    \item \textbf{Champion Model}: [Model type] on [Approach] data achieves [R² value] accuracy, suitable for [application]
\end{enumerate}

\conclusionbox{\textbf{Overall Conclusion:} The project demonstrates that domain-aware recovery and careful outlier treatment can preserve training signal at scale, enabling reliable global pricing models with consistent evaluation across preprocessing strategies.}

\subsection{Deployment Recommendations}

\begin{itemize}
    \item \textbf{Model Selection}: Use Approach 3 with tree-based model (XGBoost/LightGBM) for robustness and interpretability in original currency units
    \item \textbf{Feature Collection}: Prioritize area measurement; rooms/bathrooms secondary; country and property type essential for accurate predictions
    \item \textbf{Prediction Intervals}: Account for MAE of \$[TBD] when communicating price estimates to stakeholders
    \item \textbf{Monitoring}: Track prediction accuracy over time; retrain quarterly as market conditions evolve
\end{itemize}

\subsection{Future Work}

\begin{enumerate}
    \item Develop tier-specific models for each price segment with higher accuracy
    \item Incorporate temporal features (market trends, seasonal patterns)
    \item Add property-specific features (amenities, condition, proximity to transit)
    \item Implement ensemble combining multiple approaches for improved robustness
    \item Deploy model via API for real-time price prediction
\end{enumerate}

\end{document}
